\section{Enclosing area in the four-dimensional equality problem}
\label{q4:sec:enclosures}
\label{q4:enc:section}

We apply the geometric results of Section~\ref{q:geometry:section}
to a smooth connected orientable four-manifold $(X,g)$ with nonempty
compact smooth boundary $B$, complete with $B$ included, and exactly
one asymptotically Euclidean end with compact complement. Assume
$g-\delta=O_2(r^{-q})$ for $q>1$. The full-cut convention is
Definition~\ref{q4:def:cuts}, and
\[
 a_g(B)=\inf_{\Gamma}\operatorname{Area}_g(\Gamma).
\]
The common geometric hypotheses require only $q>0$ and are therefore
satisfied. They require neither a constraint equation nor
outermostness. In particular, we do not specialize a strong-decay
equality theorem to these weaker data.

\subsection{The full-perimeter realization}
\label{q4:enc:obstacle-setup}

Use the compact filling $O$ from Section~\ref{q:geometry:section}.
As in Equation~\eqref{q:geometry:full-perimeter}, minimize full
perimeter among bounded filled sets containing $O$; write
$\mathcal A$ for this class.

\begin{proposition}[Full-perimeter minimizing enclosures]
\label{q4:enc:minimizing-hull}
The minimum of $P_g$ over $\mathcal A$ exists and equals $a_g(B)>0$.
Every minimizer has a regular representative whose frontier $T$ is a
nonempty compact embedded $C^{1,1}$ hypersurface, smooth and minimal
on $T\setminus B$. Its graph functions are locally $W^{2,\infty}$,
hence $W^{2,2}$, including at contact with $B$. Its outward
coorientation points toward a connected exterior containing the
asymptotic end. All minimizing frontiers lie in one compact subset
of $X$.

Every frontier is a $C^1$ and area limit of smooth full enclosing
cuts in $\operatorname{int}X$. If $H_B<0$ for the normal into $X$,
every minimizer contains a fixed exterior collar of $B$. Its frontier
is then a smooth minimal full cut, and its closed exterior is smooth,
connected, complete with its nonempty boundary included, and
one-ended. Its boundary is outer area-minimizing there, with every
component counted.
\end{proposition}

\begin{proof}
Apply Proposition~\ref{q:geometry:enclosure} with $n=4$ and obstacle
$B$. Its $C^{1,1}$ conclusion applies to the whole frontier in this
dimension. Its outward collar isotopy supplies the smooth-cut
approximation, preserving the connected exterior. If $H_B<0$, the
collar flux inequality~\eqref{q:geometry:detachment} removes all
contact uniformly. The final part of the same proposition proves
intrinsic completeness and outer area-minimization of the detached
closed exterior.
\end{proof}

\paragraph{An equivalent filled-side realization.}
The argument above also admits a noncompact fixed side in place of
the compact filling.
Let $X\cup_B O'$ be the topological double, with $O'$ the second copy
of $X$, and choose a smooth metric extension through $B$ agreeing with
that of $\widehat X$ in a collar. Its admissible sets $E'$ contain
$O'$ almost everywhere, have locally finite perimeter with finite
total perimeter, and satisfy
$E'\cap\operatorname{int}X\subset K$ almost everywhere for some compact
$K\subset X$. No finite-volume condition is imposed on $E'$.

Write $P$ for full perimeter in the respective ambient manifold and
$F=E\cap\operatorname{int}X$. The maps
\[
 E=O\cup F\longmapsto E'=O'\cup F,
 \qquad E'\longmapsto O\cup(E'\cap\operatorname{int}X)
\]
are inverse modulo null sets: both retain the same essentially
compactly supported exterior part $F$. To verify that they preserve
full perimeter, let $t_F=\operatorname{Tr}_X\mathbf1_F$ be the BV
trace on $B$ taken from $X$. The gluing formula gives
\[
 P(E)=P_g(F;\operatorname{int}X)
       +\int_B|1-t_F|\,dA_g=P(E').
\]
Indeed, Gauss--Green applied on the two sides of $B$ gives the
interior variation of $\mathbf1_F$ and the jump between its trace
$t_F$ and the constant inner trace $1$. The fixed side has no
interior variation. The same calculation with compactly supported
test fields applies to $O'$, so it creates no term at infinity.
The identity also holds after restricting the perimeter measures
to any Borel subset of $X$. Thus contact with $B$ is counted in both
constructions, their perimeter-support frontiers agree, and their
minimizers correspond with the same minimizing value.

\subsection{The family of minimizers and the metric derivative}
\label{q4:enc:variation-setup}

For a minimizing filled set with frontier $T$, let $V_T$ be its
unoriented tangent-plane area measure from
Equation~\eqref{q:geometry:plane-measure}. Let $\mathcal T$ be the
family of minimizing filled sets, modulo null sets, with their
regular frontiers, equipped with local $L^1$ and weak plane-measure
convergence.

\begin{proposition}[Compact minimizing space and right area derivative]
\label{q4:enc:area-derivative}
The space $\mathcal T$ is compact and metrizable. For every smooth
symmetric two-tensor $h$, the function
\begin{equation}\label{q4:enc:metric-area-functional}
 a'_T(h)=\frac12\int_T\tr_{T_xT}h\,dA_g
\end{equation}
is continuous on it.
Let $g_s$, $|s|<s_0$, be a smooth path of smooth metrics on $X$,
with $g_0=g$, uniformly comparable to $g$, and with uniformly bounded
first and second parameter derivatives relative to $g$. Assume that
its coordinate spheres have strictly positive outward mean curvature
beyond one common radius. If $a(s)$ is its full smooth-cut infimum
and $h=\partial_sg_s|_0$, then
\begin{equation}\label{q4:enc:area-derivative-formula}
 a'(0+)=\min_{T\in\mathcal T}a'_T(h).
\end{equation}
The minimum is attained. For $s_j\to0$ and any minimizing frontiers
$T_j$ for $g_{s_j}$, a subsequence of the filled sets converges to a
member of $\mathcal T$, and its plane-area measures in the fixed
metric $g$ converge to the corresponding measure.
\end{proposition}

\begin{proof}
Smoothly extend the path through the fixed boundary in one collar.
The common mean-convex tail gives uniform confinement directly by
the clipping inequality~\eqref{q:geometry:clipping}: the divergence
of the outward unit radial normal is positive for every metric in
the path beyond the same radius. This is the common-sphere
alternative in Proposition~\ref{q:geometry:envelope}; the scaled
first-derivative bounds used in its separate density-estimate
alternative are unnecessary here. On the resulting fixed compact
region, the smooth path has uniform local bounds. The proposition's
localized polar-measure continuity and uniform Grassmann-bundle
Taylor expansion give all the assertions, retaining any mass on $B$.
\end{proof}

\begin{proposition}[A common tangent plane for tied minimizers]
\label{q4:enc:common-plane}
On $\mathcal L=\bigcup_{T\in\mathcal T}(T\setminus B)$, the plane
$\Pi_x=T_xT$ is independent of the minimizing frontier through $x$
and is continuous in the relative topology. Consequently
$x\mapsto\tr_{\Pi_x}K$ is single valued and continuous there for
every smooth symmetric two-tensor $K$.
\end{proposition}

\begin{proof}
All free frontiers are smooth in dimension four.
Proposition~\ref{q:geometry:sheets} gives local agreement at a meeting
point and continuity as both the point and the minimizing frontier
vary. Its proof uses the small-excess theorem on the whole support,
with a smooth-ambient mean-curvature bound after rescaling. This
supplies the pointwise plane control beyond the integral convergence
in Proposition~\ref{q4:enc:area-derivative}.
\end{proof}
